% --- CORRECCION: Forzar interpretacion raw de la entrada ---
\UseRawInputEncoding
\documentclass[aspectratio=169, 10pt, xcolor=table]{beamer}

% --- Configuracion del Tema ---
\usetheme{Madrid}
\usecolortheme{dolphin}

% --- Paquetes necesarios ---
\usepackage[utf8]{inputenc}
\usepackage[spanish]{babel}
\usepackage{amsmath, amssymb}
\usepackage{booktabs}
\usepackage{graphicx}
\usepackage{listings}
\usepackage{tikz}
\usetikzlibrary{arrows.meta, positioning, shapes.geometric}
\usepackage{etoolbox}

% --- Ruta de Imagenes ---
\graphicspath{{./imagenes/}}

% --- Estilo para codigo Python ---
\definecolor{codegreen}{rgb}{0,0.6,0}
\definecolor{codegray}{rgb}{0.5,0.5,0.5}
\definecolor{codepurple}{rgb}{0.58,0,0.82}
\definecolor{backcolour}{rgb}{0.95,0.95,0.92}
\lstset{
    language=Python,
    backgroundcolor=\color{backcolour},
    commentstyle=\color{codegreen},
    keywordstyle=\color{blue},
    numberstyle=\tiny\color{codegray},
    stringstyle=\color{codepurple},
    basicstyle=\ttfamily\scriptsize,
    breaklines=true,
    frame=single,
    showstringspaces=false,
    literate={á}{{\'a}}1 {é}{{\'e}}1 {í}{{\'i}}1 {ó}{{\'o}}1 {ú}{{\'u}}1 {ñ}{{\~n}}1 {ü}{{\"u}}1
}

% --- Pie de pagina personalizado ---
\makeatletter
\setbeamertemplate{footline}{
  \leavevmode%
  \hbox{%
  \begin{beamercolorbox}[wd=.333333\paperwidth,ht=2.25ex,dp=1ex,center]{author in head/foot}%
    \usebeamerfont{author in head/foot}\insertshortauthor
  \end{beamercolorbox}%
  \begin{beamercolorbox}[wd=.333333\paperwidth,ht=2.25ex,dp=1ex,center]{title in head/foot}%
    \usebeamerfont{title in head/foot}Sesion 12
  \end{beamercolorbox}%
  \begin{beamercolorbox}[wd=.333333\paperwidth,ht=2.25ex,dp=1ex,right]{date in head/foot}%
    \usebeamerfont{date in head/foot}CALCULO Y ALGEBRA LINEAL \hfill \insertframenumber/\inserttotalframenumber \hspace*{0.3cm}
  \end{beamercolorbox}}%
  \vskip0pt%
}
\makeatother

% --- Datos de la portada ---
\title[SESION 12 CALCULO]{\textbf{Sesion 12}}
\subtitle{Autovalores y Autovectores: Las Direcciones Clave}
\author[Prof. C. Sanchez]{Carlos Cesar Sanchez Coronel}
\institute{\textbf{CALCULO Y ALGEBRA LINEAL PARA IA}}
\date{}

% --- Eliminar barra de navegacion ---
\setbeamertemplate{navigation symbols}{}

% --- Indice al inicio de cada seccion ---
\AtBeginSection[]{
  \begin{frame}
    \frametitle{Contenido}
    \tableofcontents[currentsection]
  \end{frame}
}

\begin{document}

% --- Portada ---
\begin{frame}
    \titlepage
\end{frame}

% --- Indice general ---
\begin{frame}{Contenido}
    \tableofcontents
\end{frame}

% ============================================================================
% SECCION 1: DEFINICION Y SIGNIFICADO GEOMETRICO
% ============================================================================
\section{Definicion y Significado Geometrico}

\begin{frame}{Definicion de autovalor y autovector}
    \begin{block}{Definicion}
        Sea $A\in\mathbb{R}^{n\times n}$ una matriz cuadrada. Un escalar $\lambda$ y un vector no nulo $\mathbf{v}$ tales que
        \[
        A\mathbf{v} = \lambda \mathbf{v}
        \]
        se denominan \textbf{autovalor} (o valor propio) y \textbf{autovector} (o vector propio) de $A$, respectivamente.
    \end{block}
    \begin{itemize}
        \item $\mathbf{v}$ no cambia de direccion bajo la transformacion lineal.
        \item $\lambda$ es el factor de escala (puede ser positivo, negativo, cero o complejo).
    \end{itemize}
\end{frame}

\begin{frame}{Interpretacion geometrica}
    \centering
    \begin{tikzpicture}[scale=0.8]
        \draw[->] (-1,0) -- (3,0) node[right] {$x$};
        \draw[->] (0,-1) -- (0,3) node[above] {$y$};
        \draw[thick, blue, ->] (0,0) -- (1,1) node[above] {$\mathbf{v}$};
        \draw[thick, red, ->] (0,0) -- (2,2) node[above] {$A\mathbf{v}=2\mathbf{v}$};
        \draw[dashed] (1,1) -- (2,2);
    \end{tikzpicture}
    \begin{itemize}
        \item Si $\lambda>1$: el vector se alarga.
        \item Si $0<\lambda<1$: el vector se acorta.
        \item Si $\lambda<0$: el vector se invierte y se escala por $|\lambda|$.
        \item Si $\lambda=0$: el vector va al origen (nucleo).
    \end{itemize}
\end{frame}

% ============================================================================
% SECCION 2: CALCULO DE AUTOVALORES Y AUTOVECTORES
% ============================================================================
\section{Calculo de Autovalores y Autovectores}

\begin{frame}{Ecuacion caracteristica}
    \begin{block}{Ecuacion caracteristica}
        De $A\mathbf{v}=\lambda\mathbf{v}$ tenemos $(A-\lambda I)\mathbf{v}=0$. Para que exista $\mathbf{v}\neq0$, la matriz $A-\lambda I$ debe ser singular:
        \[
        \det(A-\lambda I)=0
        \]
        Este es un polinomio de grado $n$ en $\lambda$, llamado polinomio caracteristico.
    \end{block}
\end{frame}

\begin{frame}{Ejemplo manual para $2\times2$}
    \begin{exampleblock}{Ejemplo}
        Sea $A = \begin{pmatrix}4&1\\2&3\end{pmatrix}$.
        \[
        \det(A-\lambda I) = \det\begin{pmatrix}4-\lambda&1\\2&3-\lambda\end{pmatrix} = (4-\lambda)(3-\lambda)-2 = \lambda^2 -7\lambda +10 = (\lambda-5)(\lambda-2)
        \]
        Autovalores: $\lambda_1=5$, $\lambda_2=2$.
        \vspace{0.2cm}
        Para $\lambda_1=5$: $(A-5I)\mathbf{v}=0 \Rightarrow \begin{pmatrix}-1&1\\2&-2\end{pmatrix}\begin{pmatrix}x\\y\end{pmatrix}=0 \Rightarrow y=x$, autovector $\mathbf{v}_1=(1,1)^\top$.
        Para $\lambda_2=2$: $(A-2I)\mathbf{v}=0 \Rightarrow \begin{pmatrix}2&1\\2&1\end{pmatrix}\begin{pmatrix}x\\y\end{pmatrix}=0 \Rightarrow y=-2x$, autovector $\mathbf{v}_2=(1,-2)^\top$.
    \end{exampleblock}
\end{frame}

% ============================================================================
% SECCION 3: PROPIEDADES IMPORTANTES
% ============================================================================
\section{Propiedades Importantes}

\begin{frame}{Propiedades de autovalores}
    \begin{itemize}
        \item \textbf{Traza}: $\operatorname{tr}(A) = \sum_{i=1}^n A_{ii} = \sum_{i=1}^n \lambda_i$.
        \item \textbf{Determinante}: $\det(A) = \prod_{i=1}^n \lambda_i$.
        \item Si $A$ es \textbf{simetrica} ($A^\top = A$), todos los autovalores son reales y los autovectores pueden elegirse ortonormales.
        \item Si $A$ es \textbf{definida positiva}, todos sus autovalores son positivos.
    \end{itemize}
\end{frame}

% ============================================================================
% SECCION 4: DIAGONALIZACION
% ============================================================================
\section{Diagonalizacion}

\begin{frame}{Diagonalizacion de una matriz}
    \begin{block}{Diagonalizacion}
        Si $A$ tiene $n$ autovectores linealmente independientes, podemos formar la matriz $P$ con esos autovectores como columnas, y entonces
        \[
        A = P\Lambda P^{-1},\quad \Lambda = \operatorname{diag}(\lambda_1,\dots,\lambda_n)
        \]
        Esto simplifica el calculo de potencias: $A^k = P\Lambda^k P^{-1}$.
    \end{block}
    \begin{itemize}
        \item Si $A$ es simetrica, $P$ puede ser ortogonal ($P^{-1}=P^\top$).
    \end{itemize}
\end{frame}

% ============================================================================
% SECCION 5: APLICACIONES EN IA
% ============================================================================
\section{Aplicaciones en Inteligencia Artificial}

\begin{frame}{PCA (Analisis de Componentes Principales)}
    \begin{itemize}
        \item Dados datos centrados $X\in\mathbb{R}^{n\times d}$, la matriz de covarianza es $\Sigma = \frac{1}{n-1}X^\top X$.
        \item Los autovalores de $\Sigma$ son las varianzas en las direcciones de los autovectores.
        \item Seleccionamos los $k$ autovectores con mayores autovalores (componentes principales) y proyectamos: $X_{\text{proj}} = X V_k$.
        \item Reduccion de dimensionalidad conservando la maxima varianza.
    \end{itemize}
\end{frame}

\begin{frame}{PageRank de Google}
    \begin{itemize}
        \item Modela la importancia de paginas web mediante un grafo dirigido.
        \item La matriz de transicion $G$ es estocastica (columnas suman 1).
        \item El vector de ranking $\mathbf{r}$ satisface $G\mathbf{r} = \mathbf{r}$, es decir, es autovector con autovalor $\lambda=1$.
        \item Se calcula iterativamente con el metodo de la potencia.
    \end{itemize}
\end{frame}

\begin{frame}{Analisis de redes neuronales}
    \begin{itemize}
        \item En RNN, los autovalores de la matriz de pesos recurrente determinan la estabilidad del gradiente.
        \item $|\lambda|>1$ $\Rightarrow$ gradiente explosivo; $|\lambda|<1$ $\Rightarrow$ gradiente desvaneciente.
        \item El espectro del Hessiano de la perdida informa sobre la curvatura y la convergencia.
    \end{itemize}
\end{frame}

% ============================================================================
% SECCION 6: PYTHON CON NUMPY
% ============================================================================
\section{Python: Calculo y Aplicaciones}

\begin{frame}[fragile]{Calculo de autovalores y autovectores}
    \begin{lstlisting}
import numpy as np

A = np.array([[4,1],[2,3]])
eigvals, eigvecs = np.linalg.eig(A)
print("Autovalores:", eigvals)
print("Autovectores:\n", eigvecs)

# Verificacion
for i in range(len(eigvals)):
    v = eigvecs[:,i]
    lambda_i = eigvals[i]
    print(f"A @ v_{i} =", A @ v)
    print(f"lambda_{i} * v_{i} =", lambda_i * v)
    \end{lstlisting}
\end{frame}

\begin{frame}[fragile]{PCA desde cero en el dataset Iris}
    \begin{lstlisting}
import numpy as np
import matplotlib.pyplot as plt
from sklearn.datasets import load_iris
from sklearn.preprocessing import StandardScaler

iris = load_iris()
X = iris.data
scaler = StandardScaler(with_std=False)
X_centered = scaler.fit_transform(X)

cov = np.cov(X_centered.T)
eigvals, eigvecs = np.linalg.eig(cov)
idx = np.argsort(eigvals)[::-1]
eigvals = eigvals[idx]
eigvecs = eigvecs[:,idx]
print("Autovalores:", eigvals)
print("Proporcion varianza:", eigvals/np.sum(eigvals))

k=2; W = eigvecs[:,:k]; X_pca = X_centered @ W
plt.scatter(X_pca[:,0], X_pca[:,1], c=iris.target)
plt.xlabel('PC1'); plt.ylabel('PC2'); plt.show()
    \end{lstlisting}
\end{frame}

\begin{frame}[fragile]{PCA con SVD (mas estable)}
    \begin{lstlisting}
U, S, Vt = np.linalg.svd(X_centered, full_matrices=False)
# Vt contiene los autovectores (filas)
X_pca_svd = X_centered @ Vt.T[:,:2]
# Autovalores desde valores singulares
eigvals_svd = S**2 / (len(X)-1)
print("Autovalores desde SVD:", eigvals_svd)
    \end{lstlisting}
\end{frame}

\begin{frame}[fragile]{PageRank simplificado (metodo de la potencia)}
    \begin{lstlisting}
def power_iteration(A, num_iter=100):
    v = np.random.rand(A.shape[1])
    v = v / np.linalg.norm(v)
    for _ in range(num_iter):
        v = A @ v
        v = v / np.linalg.norm(v)
    eigenvalue = v.T @ A @ v
    return eigenvalue, v

# Matriz de transicion (estocastica)
A = np.array([[0.1,0.2,0.7],[0.3,0.5,0.2],[0.6,0.3,0.1]])
A = A / A.sum(axis=0, keepdims=True)
eigval, eigvec = power_iteration(A)
ranking = eigvec / eigvec.sum()
print("Ranking:", ranking)
    \end{lstlisting}
\end{frame}

% ============================================================================
% SECCION 7: NOTACION EN PAPERS
% ============================================================================
\section{Notacion en Papers Cientificos}

\begin{frame}{Notacion comun}
    \begin{itemize}
        \item $A = Q\Lambda Q^\top$: descomposicion espectral (matriz simetrica).
        \item $X^\top X \mathbf{v}_j = \lambda_j \mathbf{v}_j$: autovectores de la matriz de covarianza.
        \item $\Sigma = \sum_{i=1}^d \lambda_i \mathbf{u}_i\mathbf{u}_i^\top$: teorema espectral.
        \item $G\mathbf{r} = \mathbf{r}$: PageRank.
        \item $\nabla^2\mathcal{L}(\mathbf{w})$: Hessiano, cuyos autovalores indican curvatura.
    \end{itemize}
\end{frame}

% ============================================================================
% SECCION 8: EJERCICIOS PROPUESTOS
% ============================================================================
\section{Ejercicios Propuestos}

\begin{frame}{Ejercicios}
    \begin{enumerate}
        \item Calcula autovalores y autovectores de $A = \begin{pmatrix}2&0\\0&-3\end{pmatrix}$ e interpreta geometricamente.
        \item Demuestra que si $A$ es simetrica, sus autovalores son reales.
        \item Aplica PCA al dataset de digitos (load\_digits) y visualiza las dos primeras componentes.
        \item Investiga el metodo de la potencia y por que converge al autovector dominante.
        \item Lee el paper original de PageRank e identifica la notacion.
    \end{enumerate}
\end{frame}

% ============================================================================
% SECCION 9: CASO PRACTICO INTEGRADOR (PCA en 2D)
% ============================================================================
\section{Caso Practico: Analisis de Componentes Principales}

\begin{frame}{Problema}
    Dada la matriz de covarianza $\Sigma = \begin{pmatrix}5&2\\2&3\end{pmatrix}$, calcular:
    \begin{enumerate}
        \item Autovalores y autovectores.
        \item La primera componente principal.
        \item Porcentaje de varianza explicada.
    \end{enumerate}
\end{frame}

\begin{frame}{Solucion}
    \begin{block}{a) Ecuacion caracteristica}
        $\det\begin{pmatrix}5-\lambda&2\\2&3-\lambda\end{pmatrix} = (5-\lambda)(3-\lambda)-4 = \lambda^2-8\lambda+11=0$
        \[
        \lambda = \frac{8\pm\sqrt{64-44}}{2} = 4\pm\sqrt{5} \approx 6.236,\; 1.764
        \]
    \end{block}
    \begin{block}{b) Autovectores}
        Para $\lambda_1=4+\sqrt{5}$: $(5-\lambda_1)x+2y=0 \Rightarrow (1-\sqrt{5})x+2y=0 \Rightarrow y=\frac{\sqrt{5}-1}{2}x$.
        Un autovector: $\mathbf{v}_1 = \begin{pmatrix}2\\ \sqrt{5}-1\end{pmatrix}$ (o cualquier multiplo).
    \end{block}
    \begin{block}{c) Varianza explicada}
        $\lambda_1+\lambda_2=8$. La primera componente explica $\frac{6.236}{8}\approx 77.95\%$.
    \end{block}
\end{frame}

% ============================================================================
% SECCION 10: RESUMEN
% ============================================================================
\section{Resumen}

\begin{frame}{Lo que hemos aprendido}
    \begin{exampleblock}{Conceptos clave}
        \begin{itemize}
            \item \textbf{Autovalores y autovectores}: direcciones que no cambian bajo transformacion lineal.
            \item \textbf{Calculo}: resolviendo $\det(A-\lambda I)=0$ y luego $(A-\lambda I)\mathbf{v}=0$.
            \item \textbf{Propiedades}: traza, determinante, matrices simetricas.
            \item \textbf{Diagonalizacion}: $A = P\Lambda P^{-1}$.
            \item \textbf{Aplicaciones}: PCA, PageRank, analisis de RNN.
            \item \textbf{Python}: `np.linalg.eig`, `svd`, PCA, metodo de la potencia.
        \end{itemize}
    \end{exampleblock}
\end{frame}

\begin{frame}{Conexion con futuros cursos}
    \begin{itemize}
        \item \textbf{Machine Learning}: PCA, SVD, analisis de componentes.
        \item \textbf{Deep Learning}: espectro del Hessiano, inicializacion de pesos.
        \item \textbf{Optimizacion}: condiciones de segundo orden.
        \item \textbf{Teoria de grafos}: PageRank, centralidad espectral.
    \end{itemize}
\end{frame}

% --- Diapositiva final ---
\begin{frame}
    \centering
    \Huge \textbf{Gracias!}

\end{frame}

\end{document}